% Scaffold copied from thesis/planned-toc.md.
% This file intentionally contains headings and local TODO/context comments only.
% Old thesis prose lives in thesis/legacy/ as source material.

\section{Preliminaries}\label{sec:preliminaries}
%
% Purpose:
% - define the mathematical objects needed by later chapters.
% - show the simple, robust way to think about the domain, not hacks.
% - collapse notation and conventions that the literature isn't consistent about.
%
\subsection{Polytopes}\label{subsec:preliminaries-polytopes}
%
% Content:
%
% - polytopes in `R^4`
% - duality between convex hulls of finite sets and bounded intersections of half-spaces
% - polytopes containing zero: the dual polytope with dual vertices
% - validity checks for a proposed finite dual-row set:
%   `{a_k}` must be exactly the extremal points of `conv {a_k}`, and
%   the intersection `{x : <a_k,x> <= 1}` is bounded exactly when
%   `0 in int conv {a_k}`.
% - algorithmic construction of the primal vertex set from valid dual vertices.
%   Incidence of a primal vertex with a facet is then the defining equality
%   `<a_i,x> = 1`; precomputing incidence is useful for later algorithms but is
%   not a separate mathematical assumption.
% - support and gauge functions
% - the face lattice, names
% - Convex bodies, Convex smooth bodies, Hausdorff distance
% - The topological space of polytopes with a fixed number of facets.
% - Example: HKO2024 as a ten-facet polytope.
% - Algorithm: Volume of a polytope
%
\subsection{Smooth Symplectic Geometry}\label{subsec:preliminaries-smooth-symplectic-geometry}
%
% Content:
%
% - Standard symplectic structure and notation on `R^4`.
% - Reeb vector fields for smooth convex bodies.
% - Reeb trajectories and Reeb orbits.
% - Action of a curve.
% - Minimum action
% - Cited without proof: existence of a Reeb orbit, existence of a minimum action
% - Symplectic capacities.
% - Cited without proof: Minimum action is a symplectic capacity.
% - Viterbo's conjecture for smooth convex bodies.
% - Continuity argument via scaling and Hausdorff distance.
%
\subsection{Clarke dual action principle}\label{subsec:preliminaries-clarke-dual-action-principle}
%
% Content:
%
% - State the principle at the level needed for the thesis.
% <!-- TOC DECISION: Include the fleshed-out polished proof here. This is not just
%      motivational background. The section should name the function space,
%      action functional, minimizer/existence statement, and the conclusion used
%      later for polytopes. -->
% - Use it to justify the later finite computation story.
% <!-- OUTLINE GAP: Say which later step it justifies: existence of a
%      minimum-action orbit, reduction to a closed curve/action minimization
%      problem, or the HK2017 finite optimization problem. -->
% - Keep functional-analytic details only where they are needed for correctness.
% <!-- TOC DECISION: Because the polished proof already exists, do not hide all
%      functional-analytic details by default. Decide only how to pace the proof
%      so it remains readable. -->
%
